\scanpage{67}
% archive foliation: 66
Notons le \uline{th. de dualité absolue} :

Soient $M$ et $M^*$ deux \inserted{(complexes de)} motifs sur $X$ (de type fini sur $\mathbf{Z}$), qui sont ``duals'' pour le \deleted{$\unsure{[|\xi}$} complexe de motifs dualisants \deleted{$\unsure{]|\xi}$} canonique --- [NB Si $X$ est régulier \inserted{partout} de dim $n$, alors c'est la dualité à valeurs dans \deleted{$\unsure{\mathbf{Q}(n)[2n]|\mathbf{Q}(n)[n]}$}\inserted{$\mathbf{Q}(n)[2n+1]$}\margin{shift}\deleted{\unsure{shift}}]

alors $\mathbf{R}\Gamma_X(M)$ et $\unsure{R|\mathbf{R}}_!\Gamma_X(M^*)$ \inserted{``\unsure{support pr} sur $\mathrm{Spec}\,\mathbf{Z}$''} sont duals. Ici [Comme les coefficients sont \unsure{ess.} $\mathbf{Q}$\unsure{\}}, donc il n'y a pas de $\mathrm{Tor}$]

ceci revient aussi à la formule (cas régulier)
\[
\boxed{\begin{array}{lll}
H^i(X,M)&\text{dual de}&\mathrm{Ext}_!^{2n+1-i}(X;M,\mathbf{Q}(n))\\
H^i_!(X,M)&\text{dual de}&\mathrm{Ext}^{2n+1-i}(X;M,\mathbf{Q}(n))
\end{array}}
\]
Si $X$ propre sur $\mathrm{Spec}\,\mathbf{Z}$, inutile de mettre des $!$.

Moyennant le thm de dualité \uline{relative} pour motifs, cet énoncé revient au
